\documentclass[11pt]{article}
\usepackage[margin=1in]{geometry}
\usepackage{amsmath,amssymb,amsthm,mathtools}
\usepackage{enumitem}
\usepackage{microtype}
\usepackage[colorlinks=true,linkcolor=blue!55!black,citecolor=blue!55!black,urlcolor=blue!55!black]{hyperref}

\newtheorem{theorem}{Theorem}[section]
\newtheorem{proposition}[theorem]{Proposition}
\newtheorem{lemma}[theorem]{Lemma}
\newtheorem{corollary}[theorem]{Corollary}
\theoremstyle{remark}
\newtheorem{remark}[theorem]{Remark}

\newcommand{\Tr}{\operatorname{Tr}}
\newcommand{\dd}{\,\mathrm d}
\newcommand{\one}{\mathbf 1}

\title{Audit and Analytic Completion of the Claimed Proof\\
for the Goodman--Style Odd-Cycle Bound when $\tfrac12<p<\tfrac23$}
\author{}
\date{11 July 2026}

\begin{document}
\maketitle

\begin{abstract}
We audit the supplied claimed proof of
\[
 t(C_{2k+1},W)\ge p^{2k+1}-p(1-p)^{2k}
 \qquad \left(\tfrac12\le p\le\tfrac23\right).
\]
The operator-theoretic reduction in that proof is sound, but the submitted
version is not a satisfactory self-contained proof: one interval certificate
uses an unstated auxiliary inequality, a displayed extremum test has its
inequality reversed, and several decimal estimates are rounded in the wrong
direction.  These defects are repairable and do not produce a counterexample.
More importantly, we replace the entire computer-assisted scalar part---including
the narrow transition strip---by one analytic argument.  The new idea is to
encode the final five exponential terms as even moments of a five-atom signed
measure and prove nonnegativity of its quadratic stop-loss transform.  The only
difficult parameter regime reduces to an explicit polynomial that is positive
on a triangle.  Thus no interval subdivision, floating-point computation, or
integer-by-integer maximisation is needed in the final proof.
\end{abstract}

\section{Problem and conclusion of the audit}

Let $W\colon[0,1]^2\to[0,1]$ be a symmetric graphon and let
$p=t(K_2,W)$.  The desired bound is
\begin{equation}\label{eq:graphon-goal}
 t(C_{2k+1},W)\ge p^{2k+1}-p(1-p)^{2k}
 \qquad(k\ge1),
\end{equation}
for $1/2<p<2/3$.  Write $q=1-p$, so that $1/3<q<1/2$.

The supplied proof first performs an operator/spectral reduction.  Its final
output is the following one-parameter scalar statement.  This is the precise
interface at which the Python certificates enter.

\begin{proposition}[Scalar interface from the supplied operator reduction]
\label{prop:interface}
For every graphon in the problem and every $k\ge1$, the reduction in the
supplied note produces a number
\[
 q\le \alpha\le\sqrt q
\]
and numbers $L,c,a,\beta,\gamma$ defined by
\begin{equation}\label{eq:parameters}
 L=\frac{q-\alpha^2}{1-\alpha},\qquad
 c=1-\alpha-L,\qquad
 a=\alpha(1-\alpha),
\end{equation}
where $\beta,\gamma\ge0$ are determined by
\begin{equation}\label{eq:roots}
 \beta-\gamma=c,\qquad \beta\gamma=a,
\end{equation}
such that
\begin{equation}\label{eq:operator-reduction}
 t(C_{2k+1},W)
 \ge \beta^{2k+1}-\gamma^{2k+1}-cL^{2k}.
\end{equation}
\end{proposition}

We found no defect in the argument leading to Proposition~\ref{prop:interface}.
The remaining task is therefore to prove
\begin{equation}\label{eq:scalar-goal}
 \beta^{2k+1}-\gamma^{2k+1}-cL^{2k}
 \ge p^{2k+1}-pq^{2k}.
\end{equation}
Claude's programs certify versions of \eqref{eq:scalar-goal} by subdividing the
$(q,\alpha)$-region.  The proof below establishes \eqref{eq:scalar-goal}
analytically on the whole region at once.

\paragraph{Specific repairs to the supplied write-up.}
The original write-up should not be used unchanged.  In particular:
\begin{enumerate}[label=(\roman*),leftmargin=2.1em]
\item In the part called Region II/Zone C, the subdivision estimate uses a
lower bound on $f=\alpha-L$ that is not stated among the hypotheses.  With
$d=\alpha-q$ and $\delta=\alpha-\tfrac13$, the needed fact follows from
$d+L\le\tfrac13$:
\begin{equation}\label{eq:missing-bound}
 q-L=\alpha-(d+L)\ge\alpha-\tfrac13=\delta,
 \qquad
 f=d+(q-L)\ge d+\delta.
\end{equation}
This is stronger than the bound used by the program.
\item The displayed test for which integer maximises the auxiliary quantity
$K(n)$ has one comparison sign reversed.  The implementation and the later
use have the correct orientation, so this is a textual error rather than a
counterexample to the certificate.
\item Several handwritten decimal inequalities round an endpoint in the
unsafe direction.  Replacing those decimals by their exact rational endpoints
repairs the statements; the relevant inequalities have slack.
\end{enumerate}
These repairs explain why the computational proof is salvageable.  They are no
longer needed for the final result, since Sections~\ref{sec:scalar}--\ref{sec:proof}
replace all three certificates.

\section{The analytic scalar theorem}\label{sec:scalar}

\begin{theorem}[Five-atom scalar inequality]\label{thm:scalar}
Let
\[
 \frac13\le q\le\frac12,\qquad p=1-q,
 \qquad q\le\alpha\le\sqrt q,
\]
and define $L,c,a,\beta,\gamma$ by
\eqref{eq:parameters}--\eqref{eq:roots}.  Then, for every integer $k\ge1$,
\begin{equation}\label{eq:scalar-theorem}
 \beta^{2k+1}-\gamma^{2k+1}-cL^{2k}
 \ge p^{2k+1}-pq^{2k}.
\end{equation}
\end{theorem}

The proof uses the more convenient parameter
\begin{equation}\label{eq:fdef}
 f=\alpha-L,
 \qquad r=p-q=1-2q.
\end{equation}
Since $0\le L\le q\le\alpha$, we have $f\ge0$.  The relation
$q=\alpha^2+(1-\alpha)L$ gives the identities
\begin{equation}\label{eq:fparam}
 \alpha=\frac{q+f}{1+f},\qquad
 L=\frac{q-f^2}{1+f},\qquad
 c=\frac{r+f^2}{1+f},\qquad
 a=\frac{p(q+f)}{(1+f)^2}.
\end{equation}
In particular,
\begin{equation}\label{eq:cminus-r}
 c-r=\frac{f(f-r)}{1+f}.
\end{equation}

\begin{lemma}[Ordering of the five atoms]\label{lem:order}
The parameters satisfy
\begin{equation}\label{eq:order}
 0\le L\le q<p\le\beta,
 \qquad L\le\gamma\le p.
\end{equation}
Thus the only undetermined comparison among the interior atoms is whether
$\gamma\le q$ or $q\le\gamma$.
\end{lemma}

\begin{proof}
The bounds $0\le L\le q$ follow respectively from
$\alpha^2\le q$ and $\alpha\ge q$.  Let
\[
 \chi(x)=x^2-cx-a=(x-\beta)(x+\gamma).
\]
A direct use of \eqref{eq:fparam} gives
\begin{equation}\label{eq:chip}
 \chi(p)=-af^2\le0,
\end{equation}
so $p\le\beta$.  Moreover $a\le1/4\le p^2$, and hence
$\gamma=a/\beta\le a/p\le p$.  Finally,
\begin{equation}\label{eq:chiL}
 \chi(-L)=L^2+cL-a=-f(1-\alpha)\le0,
\end{equation}
which implies $L\le\gamma$.
\end{proof}

\section{Quadratic stop-loss domination}\label{sec:proof}

Define a signed measure on $[0,\infty)$ by
\begin{equation}\label{eq:measure}
 \nu=
 \beta\,\delta_\beta+p\,\delta_q-p\,\delta_p
 -\gamma\,\delta_\gamma-c\,\delta_L.
\end{equation}
Its total mass is zero because $\beta-\gamma=c$.  For every $k\ge1$,
\begin{equation}\label{eq:moment-is-goal}
 \int x^{2k}\,\dd\nu(x)
 =\beta^{2k+1}+pq^{2k}-p^{2k+1}
  -\gamma^{2k+1}-cL^{2k}.
\end{equation}
Consequently Theorem~\ref{thm:scalar} is equivalent to nonnegativity of all
positive even moments of $\nu$.

Consider its quadratic stop-loss transform
\begin{equation}\label{eq:Hdef}
 H(t)=\int (x-t)_+^2\,\dd\nu(x),\qquad t\ge0.
\end{equation}
The central claim is
\begin{equation}\label{eq:Hpositive}
 H(t)\ge0\qquad(t\ge0).
\end{equation}
First observe that direct substitution in \eqref{eq:fparam} gives
\begin{align}
 H(0)
 &=\beta^3+pq^2-p^3-\gamma^3-cL^2 \notag\\
 &=c(c^2+3a-L^2)-pr
 =2af^2\ge0.                                      \label{eq:H0}
\end{align}
We now prove \eqref{eq:Hpositive}.  Put
\begin{equation}\label{eq:h-tstar}
 h=\beta+\gamma,
 \qquad
 t_*=h-\frac{pr}{c}.
\end{equation}
(The open range $q<1/2$ has $c>0$; the endpoint follows by continuity.)
On every interval where the atoms at $L$ and below are inactive, the relevant
quadratic has derivative
\begin{equation}\label{eq:Hprime-core}
 2c(t-t_*).
\end{equation}
We split according to the only two support orderings and, in the second one,
according to whether $c\ge r$.

\subsection*{Case 1: $\gamma\le q$}
Here $c=\beta-\gamma\ge p-q=r$.  The intervals above $q$ are immediate from
$\beta\ge p$.  For $\gamma\le t\le q$,
\[
 H'(t)=-2\{\beta(\beta-t)-pr\}\le0,
\]
because $\beta-t\ge\beta-q\ge r$ and $\beta\ge p$.
At the lower endpoint,
\begin{align*}
 H(\gamma)
 &=\beta c^2+p(q-\gamma)^2-p(p-\gamma)^2\\
 &\ge p(q-\gamma)^2\ge0,
\end{align*}
since $c=\beta-\gamma\ge p-\gamma$ and $\beta\ge p$.
For $L\le t\le\gamma$, the derivative is \eqref{eq:Hprime-core}; moreover
$pr/c\le p\le\beta$, so $t_*\ge\gamma$.  Hence this quadratic decreases up
to $\gamma$ and is nonnegative throughout $[L,\gamma]$.

For $0\le t\le L$, all five atoms are active.  Since $\nu$ has total mass
zero, $H$ is affine there.  Its endpoint values $H(0)$ and $H(L)$ are
nonnegative, so $H\ge0$ also on this last interval.

\subsection*{Case 2: $q\le\gamma$ and $c\ge r$}
Write $y=\gamma-q\ge0$ and set
\[
 g(z)=(q+z)z^2\qquad(z\ge0).
\]
At $t=q$ we have
\begin{equation}\label{eq:Hq-superadd}
 H(q)=g(y+c)-g(y)-g(r)\ge0.
\end{equation}
Indeed, $g$ is increasing and
\[
 g(x+y)-g(x)-g(y)=xy(2q+3x+3y)\ge0;
\]
therefore $c\ge r$ implies
$g(y+c)\ge g(y+r)\ge g(y)+g(r)$.

On $[q,\gamma]$, the function $H$ is concave.  Its value at $q$ is
nonnegative, while
\[
 H(\gamma)=\beta c^2-p(p-\gamma)^2\ge0.
\]
Thus it is nonnegative on that interval.  On $[L,q]$ its derivative is
\eqref{eq:Hprime-core}.  Since
\[
 h-q\ge\beta\ge p\ge\frac{pr}{c},
\]
we have $t_*\ge q$, and hence $H(t)\ge H(q)$ on $[L,q]$.
The affine argument on $[0,L]$ finishes this case.

\subsection*{Case 3: $q\le\gamma$ and $c<r$}
This is precisely the narrow transition regime for which the supplied proof
uses its most delicate computer certificate.  By \eqref{eq:cminus-r},
\begin{equation}\label{eq:f-transition}
 0<f<r\le\frac13.
\end{equation}
Define
\begin{equation}\label{eq:u-eps}
 u=r-c=\frac{f(r-f)}{1+f}=f(c-f),
 \qquad
 \varepsilon=\beta-p.
\end{equation}
Then
\begin{equation}\label{eq:beta-gamma-eps}
 \beta=p+\varepsilon,
 \qquad
 \gamma=q+u+\varepsilon.
\end{equation}
The identity \eqref{eq:chip} is equivalent to
\begin{equation}\label{eq:eps-identity}
 \varepsilon(1+u+\varepsilon)=af^2.
\end{equation}

For $L\le t\le q$, write
\begin{equation}\label{eq:Hcore}
 H_0(t)=
 \beta(\beta-t)^2+p(q-t)^2-p(p-t)^2-\gamma(\gamma-t)^2.
\end{equation}
For $z=q-t$ this becomes
\begin{equation}\label{eq:quadratic-z}
 H_0(q-z)=H(q)+2Tz+cz^2,
\end{equation}
where direct expansion gives
\begin{align}
 T&=2c\varepsilon-u(q+u),                                      \label{eq:T}\\
 H(q)&=3c\varepsilon^2
       +c\{2q+3(r+u)\}\varepsilon-u^2(q+u).                    \label{eq:Hq}
\end{align}
It is enough to prove that the quadratic in \eqref{eq:quadratic-z} has
nonpositive discriminant.  Using \eqref{eq:eps-identity}, one obtains the
exact identity
\begin{equation}\label{eq:Didentity}
 D:=cH(q)-T^2
 =4prc\,\varepsilon-c^2af^2-pu^2(q+u).
\end{equation}

The transition condition $f<r$ implies
$\alpha=(q+f)/(1+f)\le1/2$.  Hence $a=\alpha(1-\alpha)\ge pq$.
Also $\gamma\le p$ by Lemma~\ref{lem:order}.  From
\eqref{eq:eps-identity} we therefore get
\begin{equation}\label{eq:eps-lower}
 \varepsilon=\frac{af^2}{p+\gamma}
 \ge\frac{qf^2}{2}.
\end{equation}
Since $u=f(c-f)$, equations \eqref{eq:Didentity} and
\eqref{eq:eps-lower} yield
\begin{equation}\label{eq:Dlower}
 D\ge f^2\mathcal E,
 \qquad
 \mathcal E=2pqrc-c^2a-p(c-f)^2(q+u).
\end{equation}
It remains to show $\mathcal E\ge0$.

Set
\begin{equation}\label{eq:Fdef}
 F=\frac{(1+f)^4}{p}\,\mathcal E.
\end{equation}
Substituting \eqref{eq:fparam} and \eqref{eq:u-eps} and collecting terms gives
the exact factorisation
\begin{equation}\label{eq:Pfactor}
 2F=fP(r,f),
\end{equation}
where
\begin{align}
 P(r,f)={}&2r(1-3r^2)
 +(-1+5r+7r^2-7r^3)f 
otag\\
 &-2(1+r^3)f^2
 +2r(1-3r)f^3+2r(1-r)f^4.                    \label{eq:Ppoly}
\end{align}
We prove $P(r,f)\ge0$ on the triangle $0\le f\le r\le1/3$.

If $0\le r\le1/5$, the coefficient of $f$ is at least $-1$ and the
coefficients of $f^3,f^4$ are nonnegative.  Therefore
\begin{align*}
 P(r,f)
 &\ge2r(1-3r^2)-f-2(1+r^3)f^2\\
 &\ge r(1-2r-6r^2-2r^4)\\
 &\ge \frac{223}{625}\,r\ge0.
\end{align*}
For $1/5\le r\le1/3$, the coefficient
$-1+5r+7r^2-7r^3$ is nonnegative, and hence
\begin{align*}
 P(r,f)
 &\ge2r(1-3r^2)-2(1+r^3)f^2\\
 &\ge2r(1-r-3r^2-r^4)\\
 &\ge\frac{52}{81}\,r\ge0.
\end{align*}
(The two bracketed polynomials are decreasing on their respective intervals.)
Thus $F\ge0$, so $\mathcal E\ge0$, and consequently $D\ge0$.

It follows from \eqref{eq:quadratic-z} that
\[
 4(T^2-cH(q))=-4D\le0.
\]
Hence $H_0(t)\ge0$ for every real $t$.  On $[q,\gamma]$, the function $H$ is
concave and both endpoint values are nonnegative: $H(q)\ge0$ follows from
$D\ge0$, and $H(\gamma)\ge0$ was proved above.  The intervals above $\gamma$
are immediate.  Finally, on $[0,L]$ the full $H$ is affine, with
$H(0)\ge0$ by \eqref{eq:H0} and $H(L)=H_0(L)\ge0$.
This completes the proof of \eqref{eq:Hpositive} in the transition case and
therefore in all cases.

\subsection*{From stop-loss domination to all odd cycle lengths}
For every real $s\ge3$ and $x\ge0$,
\begin{equation}\label{eq:integral-representation}
 x^s=\frac{s(s-1)(s-2)}2
 \int_0^\infty t^{s-3}(x-t)_+^2\,\dd t.
\end{equation}
Consequently \eqref{eq:Hpositive} implies
$\int x^s\,\dd\nu(x)\ge0$ for every $s\ge3$.  For the remaining even exponent
$s=2$, equation \eqref{eq:H0} gives the same conclusion directly.  Taking
$s=2k$ in \eqref{eq:moment-is-goal} proves
Theorem~\ref{thm:scalar}.

\section{Consequence for the graphon problem}

Combining Proposition~\ref{prop:interface} with
Theorem~\ref{thm:scalar} gives immediately:

\begin{corollary}[Goodman--style bound in the missing density range]
\label{cor:main}
Let $W$ be a graphon with edge density $p\in[1/2,2/3]$.  Then, for every
integer $k\ge1$,
\[
 t(C_{2k+1},W)\ge p^{2k+1}-p(1-p)^{2k}.
\]
\end{corollary}

\begin{remark}[Logical role of the checker]
The accompanying Python file checks the polynomial identities and stress-tests
$H(t)\ge0$ at high precision.  It is a regression checker only.  No line of
the proof of Theorem~\ref{thm:scalar}, and hence no part of the final scalar
argument, depends on numerical sampling or interval arithmetic.
\end{remark}

\section{What should replace the supplied computational sections}

For a revised research note, we recommend retaining the operator reduction up
to Proposition~\ref{prop:interface}, deleting the zone decomposition and all
three computer-assisted scalar certificates, and inserting
Sections~\ref{sec:scalar}--\ref{sec:proof}.  This makes the central argument
shorter in logical structure, removes the transition-strip fragility, and
proves all odd lengths simultaneously rather than verifying a succession of
auxiliary exponential comparisons.

\end{document}
